\chapter{\IfLanguageName{dutch}{Short list}{Short list}}%
\label{ch:shortlist}

Op basis van de vergelijkingstabellen in de long list worden in dit hoofdstuk de drie geselecteerde alternatieven per variabel component van de \gls{rag}-pipeline diepgaand besproken. De short list dient niet louter als bevestiging van de tabelscore, maar als een inhoudelijke onderbouwing: per component wordt aangetoond waarom precies deze drie opties de beste keuze zijn voor de specifieke context van dit onderzoek en de \gls{poc} van Turtle Srl. De combinatie van de negen geselecteerde componenten levert de reeds beschreven testmatrix van $3^3 = 27$ unieke \gls{rag}-configuraties op.

\section{Chunking-strategieën}
\label{sec:shortlist-chunking}

\subsection{Fixed-size chunking}

Fixed-size chunking wordt opgenomen als onmisbare \textbf{baseline}. Zonder een deterministisch referentiepunt is het wetenschappelijk onmogelijk om de meerwaarde van de twee complexere strategieën aan te tonen: elke benchmark heeft een nulpunt nodig. \textcite{Lewis2020} bevestigen dit door al in het originele \gls{rag}-paper chunks van vaste lengte te gebruiken als standaard retrieval-eenheid. De strategie is volledig deterministisch en vereist geen extra modelaanroepen, wat directe reproduceerbaarheid en maximale snelheid tijdens de indexeringsfase garandeert — twee Must-Have-requirements uit de requirementsanalyse.

Concreet betekent dit voor de experimentele opzet dat de chunk-grootte en overlap als configureerbare parameters worden vastgelegd via HyPSTER. \textcite{BhatEtAl2025} stellen op basis van een systematische evaluatie over zes datasets dat 512 tokens met 10\% overlap een sterke baseline vormt voor \gls{esg}-gerelateerde\\vraag-en-antwoordtaken. Deze parameterwaarden bieden een evenwicht tussen context per chunk en granulariteit van de retrieval. De beperking — het negeren van tekststructuur — is bewust geaccepteerd, omdat het net de lacune is die de twee andere strategieën moeten opvullen en meetbaar moeten verbeteren.

\subsection{Recursive character chunking}

Recursive character chunking wordt geselecteerd als de \textbf{structuurbewuste middenweg}. De strategie splitst tekst hiërarchisch op een geordende reeks separatoren: eerst alinea-grenzen (\texttt{\textbackslash n\textbackslash n}), dan regelafstanden (\texttt{\textbackslash n}), dan spaties en ten slotte losse tekens \autocite{LangChainDocs2024}. Dit respecteert de natuurlijke opbouw van \gls{esg}-rapporten, die typisch gestructureerd zijn in alinea's en bullet-lijsten per \gls{esg}-criterium, zonder daarvoor afhankelijk te zijn van metadata of embedding-aanroepen.

De empirische onderbouwing voor deze keuze is sterk. \textcite{Amiri2025} evalueerden 25 chunking-configuraties over vijf methodefamilies en 48 embedding-modellen in een domeinspecifieke context en stelden vast dat recursive token-based chunking consistent het best presteerde, met tot \textbf{45\% hogere domain-weighted precision} dan de beste fixed-size variant. Voor \gls{esg}-documenten is dit bijzonder relevant: een volledige \gls{vsme}-richtlijnpassage of een GRI-standaard-alinea wordt als één coherente chunk opgehaald in plaats van willekeurig afgekapt te worden. \textcite{SmithTroynikov2024} bevestigen in hun evaluatiestudie dat de keuze van chunking-strategie tot 9\% verschil in recall veroorzaakt, waarbij recursive character splitting in de praktijk consistent goed presteert.

Vanuit de requirementsanalyse scoort deze strategie het hoogst op de Must-Have-requirements: zij maximaliseert Context Recall door alinea-integriteit te bewaren, garandeert source attribution omdat metadata per chunk intact blijft, en is volledig implementeerbaar als Haystack \texttt{@component}. Tegelijkertijd is zij even snel en deterministisch als fixed-size chunking, wat de reproduceerbaarheid-requirement waarborgt.

\subsection{Semantic chunking}

Semantic chunking wordt geselecteerd als de \textbf{semantisch meest geavanceerde} strategie, geformaliseerd door \textcite{Kamradt2023}. De methode berekent de cosine-gelijkenis tussen opeenvolgende zinnen via hun embeddings en plaatst een chunk-grens wanneer de semantische gelijkenis onder een instelbare drempelwaarde daalt. Dit maakt het de enige strategie die actief \textit{betekenisgrenzen} detecteert in plaats van structuur- of tekengrenzen.

Voor de specifieke documentbasis van \gls{esg}-rapporten is dit bijzonder relevant: één sectie kan meerdere onderwerpen behandelen, zoals energieverbruik en \gls{co2}-uitstoot in dezelfde alinea, terwijl twee thematisch verwante alinea's door een technische kopstructuur worden gescheiden. \textcite{MorenoCediel2026} tonen aan dat semantic chunking chunks produceert met een hogere interne semantische coherentie, wat de context precision verhoogt. In een \gls{rag}-pipeline betekent dit dat de opgehaalde chunks meer gerichte informatie bevatten en minder ruis, wat de taak van de \gls{llm} vereenvoudigt en de kans op gefabriceerde aanvullingen vermindert.

Een cruciaal implementatiedetail is de drempelwaarde voor de cosine-gelijkenissplit. \textcite{Kamradt2023} stelt de drempel standaard in op het 95e percentiel van de afstandsverdeling, maar dit is documentafhankelijk. Voor de experimentele opzet van dit onderzoek wordt de drempel handmatig bepaald en gedocumenteerd per run, in lijn met de white-box requirement uit de requirementsanalyse. De hogere indexeringskosten — één embedding-aanroep per zin — zijn bewust aanvaard, omdat het onderzoek precies tot doel heeft te meten of deze meerkost wordt gerechtvaardigd door betere retrieval-kwaliteit. \textcite{QuEtAl2024} waarschuwen dat dit niet altijd het geval is, maar hun bevinding benadrukt net de wetenschappelijke waarde van het systematisch testen van alle 27 configuraties.

\section{Embedding-modellen}
\label{sec:shortlist-embedding}

\subsection{Snowflake arctic-embed-l-v2.0}

Arctic-embed-l-v2.0 wordt geselecteerd als de \textbf{meertalige open-source baseline} voor de embedding-component. De kernmotivatie voor deze selectie is tweeledig: het model scoort het sterkst van alle onderzochte open-source kandidaten op de meertalige retrieval-benchmark die het meest relevant is voor de documentbasis van Turtle Srl, en het doet dit onder een Apache 2.0-licentie die volledig vrije commerciële deployment toelaat.

De doorslaggevende prestatie-indicator is de score op de CLEF meertalige retrieval-benchmark, die nDCG@10 meet over Engels, Frans, Spaans, Italiaans en Duits. Arctic-embed-l-v2.0 behaalt een gemiddelde van 0,541 op deze benchmark, terwijl BGE-M3 uitkomt op 0,410 en multilingual-e5-large op 0,431 \autocite{YuEtAl2024}. Voor de Italiaans-Engelse documentbasis van Turtle Srl is Italiaanse retrieval-kwaliteit een harde eis — de copromotor stelde expliciet dat de pipeline meertalige documenten moet kunnen verwerken — en arctic-embed-l-v2.0 is het enige open-source model in de short list dat op dit specifieke taalprofi\-el empirisch superieur is aan de twee andere kandidaten.

Vanuit de requirementsanalyse vult dit model de Must-Have-requirements voor meertaligheid, open-source licentie en Azure-compatibiliteit volledig in. Het model integreert rechtstreeks via \texttt{SentenceTransformersDocumentEmbedder} in Haystack 2.x en heeft een maximale inputlengte van 8.192 tokens. Als baseline in het embedding-component toont het de vooruitgang met uitsluitend dense semantisch search, zonder hybride functionaliteiten.

\subsection{BAAI BGE-M3}

BGE-M3 wordt geselecteerd als het \textbf{technologisch meest veelzijdige} open-source alternatief. De unieke eigenschap die dit model onderscheidt van alle andere kandidaten in de long list is de gelijktijdige ondersteuning voor drie retrievalfunctionaliteiten via één modelgewicht \autocite{Chen2024}: dense retrieval voor semantische gelijkenis, sparse retrieval voor trefwoordmatching (vergelijkbaar met BM25) en multi-vector retrieval voor token-niveau matching in ColBERT-stijl. Dit is het eerste embedding-model dat deze drie functionaliteiten combineert via een \textit{self-knowledge distillation}-trainingstechniek.

Voor de \gls{esg}-use case van Turtle Srl is hybride zoeken bijzonder waardevol. \gls{esg}-vragenlijsten combineren conceptuele vragen (``Welke maatregelen neemt het bedrijf voor klimaatmitigatie?'') met specifieke trefwoordzoekopdrachten (``Scope 3 emissies'', ``\gls{vsme} artikel 4''). Puur semantisch zoeken mist soms exacte terminologie; puur trefwoordzoeken mist contextueel verwante passages. BGE-M3 combineert beide via \gls{rrf} in Qdrant, wat de verwachting schept van hogere Context Recall dan text-embedding-3-large \autocite{QdrantHybrid2025}. De ondersteuning voor meer dan 100 talen en een maximale inputlengte van 8.192 tokens voldoen ruimschoots aan de meertaligheids- en chunk-size-eisen. Als open-source model biedt het bovendien kostefficiëntie op schaal en de mogelijkheid tot lokale deployment, wat \gls{gdpr}-gevoeligheid maximaliseert.

\subsection{Microsoft multilingual-e5-large-instruct}

Multilingual-e5-large-instruct wordt geselecteerd als het \textbf{instruction-geoptimaliseerde} alternatief. De E5-modelfamilie was de eerste die het concept van asynchrone \textit{query-} en \textit{passage-}prefixen introduceerde als expliciete instruction-tuning voor retrieval \autocite{WangEtAl2022}. De instructie-meganisme werkt als volgt: door \texttt{query:} voor een vraag en \texttt{passage:} voor een documentchunk te plaatsen, dwingt men het model de fundamenteel asymmetrische relatie tussen een vraag en zijn antwoord te modelleren \autocite{Wang2024a}. Dit sluit direct aan bij de eis van mijn copromotor voor modellen die geoptimaliseerd zijn voor instructies en sterke retrieval op basis van vragen en trefwoorden.

De multilingual variant, beschreven door \textcite{WangEtAl2024b}, initialiseert vanuit xlm-roberta-large en is getraind op 1 miljard meertalige tekstparen via contrastief pre-trainen gevolgd door instruction-tuning, met state-of-the-art resultaten op\\\gls{beir} en \gls{mteb} benchmarks voor meertalige retrieval. Een bijkomend praktisch voordeel is het lagere gewicht vergeleken met BGE-M3, wat snellere lokale inferentie op een Azure-server mogelijk maakt en de latency-requirement ondersteunt. De beperking is de afwezigheid van hybride zoekfunctionaliteit, waardoor het model uitsluitend op dense semantische vectors vertrouwt.

\section{Large Language Models}
\label{sec:shortlist-llm}

\subsection{Google Gemma 3 27B IT}

Gemma 3 27B IT wordt geselecteerd als de \textbf{lichtgewicht instruction-geoptimaliseerde open-source baseline} voor de generatieve component van de pipeline. De keuze voor dit model als baseline combineert twee argumenten: de faithfulness-prestaties zijn sterk genoeg om als referentiepunt te dienen, en de inferentiesnelheid is bij 27 miljard parameters hoog genoeg om de kostenefficiëntie-requirement te verankeren.

Vanuit de requirementsanalyse zijn drie eigenschappen doorslaggevend. Ten eerste is de \textbf{instruction-tuning} expliciet ingebakken: het achtervoesel \textit{IT} staat voor \textit{Instruction Tuned} — Gemma 3 27B IT is post-training gefinetuned voor het opvolgen van complexe systeemprompts en het genereren van gestructureerde \gls{json}-output, wat cruciaal is voor de Metadata-Aware Prompt Builder in de Haystack-pipeline \autocite{Gemma3Team2025}. Ten tweede zijn de meertalige prestaties sterk: het model werd getraind op een divers meertalig corpus en ondersteunt Italiaans en Engels, wat de meertaligheidseis van Turtle Srl invult. Ten derde is het model volledig \textbf{\gls{gdpr}-compliant} inzetbaar via lokale deployment op een Azure VM of via de Azure AI Foundry Model Catalog — geen data verlaat de EU-infrastructuur. De Gemma-licentie staat commercieel gebruik toe, wat de open-source requirement invult \autocite{Gemma3Team2025}. Met 27 miljard parameters draait het model vlot op een enkele Azure VM, wat de latency-requirement ondersteunt. Als baseline vormt Gemma 3 27B IT het referentiepunt voor de efficiëntiedimensie van de benchmark: het toont hoe goed een compact model presteert en of de extra infrastructuurkost van Llama 3.3 70B daadwerkelijk betere \gls{ragas}-scores rechtvaardigt.

\subsection{Meta Llama 3.3 70B Instruct}

Llama 3.3 70B Instruct wordt geselecteerd als het \textbf{open-source model met de hoogste faithfulness en instructieopvolgingskwaliteit}. De opname van dit model stelt de kernvraag van de \gls{llm}-vergelijking scherp: rechtvaardigt de extra infrastructuurkost van een 70B-model meetbaar betere \gls{ragas}-scores ten opzichte van het compactere Gemma 3 27B IT?

\textcite{GrattafioriEtAl2024} beschrijven hoe Llama 3.3 70B Instruct werd uitgebracht als een verbeterde post-getrainde versie van de Llama 3 70B-architectuur. De benchmarkcijfers positioneren het als het sterkste model in de short list op de twee meest kritische requirements. Op faithfulness haalt het een hallucinatiescore van 4,1\% op de Vectara Hallucination Leaderboard \autocite{TamberEtAl2025} — lager dan Gemma 3 27B IT en lager dan Mistral Small 4, wat het tot het meest faithfulness-betrouwbare model van de drie maakt. De IFEval-score van 92,1\% is de hoogste van alle onderzochte kandidaten. IFEval meet precies hoe betrouwbaar een model complexe prompt-instructies opvolgt, inclusief outputformatering, weigering van buiten-context-vragen en bronvermeldingsconventies. Dit is de meest directe predictor van hoe goed het model de Haystack-prompt-instructies voor de Metadata-Aware Prompt Builder zal naleven. \gls{mmlu} van 86,0\% en \gls{mmlu}-Pro van 68,9\% bevestigen de brede redeneer- en taalbegripscapaciteiten. Italiaans is een van de acht officieel ondersteunde talen, wat de meertaligheidseis volledig invult. De beschikbaarheid als eerste partij op Microsoft Azure Foundry met \gls{eu}-dataresidentie maakt Azure-deployment eenvoudig en garandeert \gls{gdpr}-naleving zonder bijkomende configuratie. De bewuste beperking is de hogere hardwarevereiste: voor vlotte inferentie is een grotere \gls{gpu}-instantie nodig dan voor Gemma 3 27B IT, wat hogere Azure-kosten meebrengt.

\subsection{Mistral AI Mistral Small 4 (mistral-small-4-119b-2603)}

Mistral Small 4 wordt geselecteerd als het \textbf{\gls{eu}-soevereine instruction-geoptimaliseerde alternatief}. De selectie van dit model introduceert precies de dimensie die Mistral Large 2 — het model dat initieel deze positie innam — niet kon bieden: gecertificeerde Europese datasoevereiniteit gecombineerd met volwaardige instruction-tuning. Mistral Large 2 is geen instruction-tuned variant en mist daarmee de post-training fine-tuning die nodig is voor betrouwbare instructieopvolging in de chatbot-context van Turtle Srl.

Mistral AI is gevestigd in Parijs en verwerkt alle data exclusief binnen de \gls{eu}, waardoor Mistral Small 4 niet valt onder de Amerikaanse CLOUD Act \autocite{MistralAISmall42026}. Dit maakt het het enige model in de short list dat absolute Europese datasoevereiniteit garandeert én instruction-geoptimaliseerd is voor chatbot-toepassingen. Het model ondersteunt meerdere talen waaronder Italiaans en Engels, beschikt over een context window van 128.000 tokens en is beschikbaar onder de Apache 2.0-licentie \autocite{MistralAISmall42026}. De Azure AI Foundry Model Catalog biedt native ondersteuning, wat de Azure-compatibiliteitseis invult, én het model ondersteunt volledig self-hosted deployment op een Azure VM via vLLM of een Docker-container — de meest flexibele productie-optie voor een toekomstige on-premise implementatie bij Turtle Srl. De hogere parametercount (119 miljard) impliceert grotere hardwarevereisten dan Gemma 3 27B IT, maar biedt tegelijkertijd sterkere redeneer- en instructieopvolgingscapaciteiten voor complexe \gls{esg}-vragen.
